% answers_short.tex —— 2004 年·全国卷Ⅰ·理科 · 简短答案速查
% 由 tools/make_short_answers.py 自 answers.tex 的 \gksolhead/\gkanswer 提取，
% 并与 2004/2004_数学真题_全国1理_答案要点.md 逐题比对后生成，勿手改（改源 MD 后重跑）。

\begin{gknotice}
\textbf{使用说明}：本节只给最终答案便于快速核对，完整推导见解析版。答案为编写者独立推导，\textbf{非官方标准答案}。
\end{gknotice}

\gkpart{选择题}
\begin{center}\begin{tabular}{ccccccccccccc}
\toprule
题号 & 1 & 2 & 3 & 4 & 5 & 6 & 7 & 8 & 9 & 10 & 11 & 12 \\
\midrule
答案 & D & B & C & B & A & B & C & C & B & A & D & B \\
\bottomrule
\end{tabular}\end{center}

\gkpart{填空题}
\begin{itemize}[leftmargin=2.6em]
  \item[\textbf{13.}] $[-1,+\infty)$
  \item[\textbf{14.}] $x^2+y^2=4$
  \item[\textbf{15.}] $a_1=1$，且 $a_n=\dfrac{n!}{2}\ (n\geqslant2)$
  \item[\textbf{16.}] ①②④
\end{itemize}

\gkpart{解答题}
\begin{itemize}[leftmargin=2.6em]
  \item[\textbf{17.}] 最小正周期为 $\pi$，最大值为 $\dfrac34$，最小值为 $\dfrac14$
  \item[\textbf{18.}] $P(\xi=0,1,2,3,4)=\left(\dfrac9{100},\dfrac3{10},\dfrac{37}{100},\dfrac15,\dfrac1{25}\right)$，且 $E\xi=\dfrac95$
  \item[\textbf{19.}] $a>0$ 时增区间为 $\left(-\infty,-\dfrac2a\right)\cup(0,+\infty)$、减区间为 $\left(-\dfrac2a,0\right)$；$a=0$ 时减区间为 $(-\infty,0)$、增区间为 $(0,+\infty)$；$a<0$ 时减区间为 $(-\infty,0)\cup\left(-\dfrac2a,+\infty\right)$、增区间为 $\left(0,-\dfrac2a\right)$
  \item[\textbf{20.}] (1) 点 $P$ 到平面 $ABCD$ 的距离为 $\dfrac32$；(2) 两平面的锐角夹角 $\theta$ 满足 $\cos\theta=\dfrac{2\sqrt7}{7}$，即 $\theta=\arccos\dfrac{2\sqrt7}{7}$；若取棱锥内部二面角，则为 $\pi-\theta$
  \item[\textbf{21.}] (1) $e\in\left(\dfrac{\sqrt6}{2},+\infty\right)\setminus\{\sqrt2\}$；(2) $a=\dfrac{17}{13}$
  \item[\textbf{22.}] (1) $a_3=3$，$a_5=13$；(2) 偶数 $n=2k$ 时 $a_n=\dfrac{3^k+(-1)^k-2}{2}$；奇数 $n=2k+1$ 时 $a_n=\dfrac{3^{k+1}+(-1)^k-2}{2}$
\end{itemize}
